\section[chaîne de numérisation]{La chaîne de numérisation de CujasNum}

La numérisation des documents anciens demandent au préalable une réflexion poussée \enquote{portant sur les choix techniques comme intellectuels}. L'avantage de l'opération de numérisation pour une bibliothèque comme celle de Cujas est l'opportunité extrêmement efficace de mettre en valeur une collection.

Malgré la richesse des collections présentent dans la réserve de la bibliothèque Cujas, celle-ci n'a pas l'effet escompté auprès des lecteurs au début. Frédérique Joannic-Seta affirme cette idée dans son artcile et la cause de ceci selon elle, est du à \enquote{ses difficultés d'ouverte de consultation qui entraînent une faible utilisation au regard de sa réelle richesse}. Par ailleurs, les ouvrages sont en mauvais états ayant subit des \enquote{inondations exceptionnelles dans les magasins que des dégradations quotidiennes suscitées par les manipulations et les conditions de cnoservation, longtemps médiocres}. 

En 1998, la volonté de la bibliothèque Cujas de pouvoir mettre en valeur cette réserve s'avère nécessaire et décide de recourir à la numérisation de ses collections. Son choix de faire de la numérisation est donc dicté par \enquote{un souci de conservation des documents par la mise à disposition des lecteurs d'un support de substitution de qualité}. Le deuxième objectif est de vouloir rendre l'accès aux documents plus facile et permettre aux lecteurs \enquote{une meilleure appropriation du texte en leur proposant des facilités de reproduction}. 

Après l'exposition virtuelle réalisée sur Jean Carbonnier, la bibliothèque Cujas possède une chaîne de numérisation en interne opérationnelle demandant à être utilisé. Celle-ci s'est renforcé gràce à la création de la célèbre \enquote{liste Pfister-Roumy}, qui est \enquote{alimentée enti-èrement en interne grâce aux compétences du département informatique}. Pour cette liste, Cujas a fait le choix de numériser en internes uniquement les \enquote{ouvrages postérieurs à 1830 avec l'océrisation et la relecture partielle sur les parties structurantes}. Pour les ouvrages antérieurs à 1830, ils sont numérisés par un partenaires, Puces et Plumes pour \enquote{préserver le document}, car le scanner que possède représente un risque pour ses ouvrages déjà fragiles. Le prestataire se charge de numériser les documents en \enquote{mode image, de faire la saisie fines des tables des matières}. 

Afin que les projets de numérisation soient bien réalisé, la méthode de travail de Cujas passe par la création d'un bordereau comportant diverses parties : la première partie correspond à la notice bibliographique du document, la seconde correspond à la pagination et aux particularités observées lors de la consultation de l'ouvrage, et la dernière partie indique l'état du document avant le transfert chez le prestataire.

Une fois, la numérisation effectuée la bibliothèque réceptionnne les images réalisées sous le format TIFF. Elle effectue alors \enquote{un contrôle physique, livre en main, pour vérifier la complétude des images et s'assurer qu'aucune page ne manque par rapport au bordereau}. Ensuite, elle génère les formats de diffusion (JPEG, PDF) et produit ou vérifie que la table de concordance \enquote{qui fait correspondre chaque numéro d'image à la pagination réelle du livre}, grâce à un script. Au moment, la personne en charge du contrôle qualité vérifie également les informations rentrées dans la Dublin Core s'assurant ainsi que tout soit correcte. Parallèlement à son travail de contrôle, des mentors étudiants sont en charge de la saisie intégrale des tables des matières. La table des matières a la particularité de représenter la structure intellectuelle de l'ouvrage. Une fois que le contrôle qualité est validé vint la dernière étape celle de la mise en publication c'est-à-dire il s'agit de vérifier \enquote{la résolution et la compression des images JPEG et TIFF dans IrfanView, de vérouiller les PDF grâce à la protection par un mot de passe}, de valider la structure XML du Dublin Core ainsi que le paquet d'archives afin d'éviter un rejet par le CINES. Cette étape terminée, les données valides sont alors \enquote{transférées via un cleint SFTP sécurisé durant la nuit vers l'espace de publication, un protocole de sécurité renforcé suite à un incident de rançonnage}. 

Au cours du traitement numérique, au démarrage la Reconnaissance Optique de Caractères (OCR) étaient appliquées uniquement aux documents postérieurs à 1830 pour permettre la recherche plein-texte. Lorsque cela n'est pas possible comme pour les ouvrages antérieurs à 1830 ceux-ci ont l'OCR s'appliquant uniquement à la table des martières, saisies entièrement à la main. 

Une autre particularité que l'on peut observer lié à la chaîne de numérisation est l'indexation que celle-ci va utiliser pour indexer ses documents numérisés dans sa bibliothèque numérique. Il s'agit de l'indexation Hauville, basée sur l'indexation RAMEAU.